\documentclass[11pt]{article}

\usepackage[margin=1.05in]{geometry}
\usepackage{amsmath,amssymb,amsthm,amsfonts}
\usepackage{booktabs}
\usepackage{microtype}
\usepackage[usenames,dvipsnames]{xcolor}
\usepackage[colorlinks=true,linkcolor=NavyBlue,citecolor=NavyBlue,urlcolor=NavyBlue]{hyperref}

\theoremstyle{definition}
\newtheorem{definition}{Definition}[section]
\theoremstyle{plain}
\newtheorem{proposition}[definition]{Proposition}
\newtheorem{theorem}[definition]{Theorem}
\theoremstyle{remark}
\newtheorem{remark}[definition]{Remark}

\newcommand{\Rb}{\mathbb{R}}
\newcommand{\Eb}{\mathbb{E}}
\newcommand{\Pb}{\mathcal{P}}
\newcommand{\KL}{\mathrm{KL}}
\newcommand{\Dt}{\mathcal{D}}
\newcommand{\Lh}{\mathcal{L}}
\newcommand{\Mcal}{\mathcal{M}}
\newcommand{\Rcal}{\mathcal{R}}
\newcommand{\proj}{\mathrm{proj}}

\title{\bfseries From Sequential Posteriors to a Schr\"odinger Bridge\\
\large A walkthrough of posterior-bridge continual learning}
\author{Working notes}
\date{\today}

\begin{document}
\maketitle

\begin{abstract}
\noindent
The method is built around three ideas, and this note develops each one only at
the point where it becomes necessary.

\S\ref{sec:problem} establishes what continual learning actually requires: not
each task's posterior but the \emph{sequential} posterior, carried forward in a
fixed-size representation because the data is discarded.
\S\ref{sec:chase} is the heart of the note --- it describes what physically
happens when we move the particles, in terms of a \emph{target that moves and
particles that chase it}, and shows where the equilibrium sits at every moment.
\S\ref{sec:lag} shows that the chase falls behind, and that this lag is the
entire reason a Schr\"odinger bridge is needed.
\S\ref{sec:reference} explains why the reference process is annealed Langevin
and not Brownian motion, which is the choice that determines what ``damage''
means in weight space.
\S\ref{sec:sb}--\S\ref{sec:imf} define the bridge and the algorithm that finds
it; \S\ref{sec:implementation}--\S\ref{sec:numbers} give the implementation and
the measurements.
\end{abstract}

\tableofcontents

% ===========================================================================
\section{What continual learning actually requires}
\label{sec:problem}

\subsection{Each task induces a posterior}

Let $\theta \in \Rb^P$ collect every weight of one shared network; for us
$P = 89{,}610$. Task $t$ arrives as a dataset $\Dt_t$ and induces
\begin{equation}
p(\theta \mid \Dt_t) \;\propto\; p(\Dt_t \mid \theta)\, p(\theta).
\end{equation}

This is a \emph{distribution}, not a point. A neural network task does not have
one solution; permuting hidden units gives identical functions, and beyond that
exact symmetry there is a wide region of weight space where accuracy is
essentially unchanged. Training returns one arbitrary member of that family.

\subsection{But we need a different object}

We do not want $p(\theta \mid \Dt_t)$. We want the \emph{sequential posterior}
\begin{equation}
q_t(\theta) \;\triangleq\; p(\theta \mid \Dt_{1:t}),
\end{equation}
which is a genuinely different distribution. Using conditional independence and
substituting $p(\Dt_i\mid\theta) \propto p(\theta\mid\Dt_i)/p(\theta)$,
\begin{equation}
p(\theta \mid \Dt_{1:t}) \;\propto\;
\frac{\prod_{i=1}^{t} p(\theta \mid \Dt_i)}{p(\theta)^{\,t-1}} .
\label{eq:poe}
\end{equation}
A \emph{product} --- the object every Bayesian continual-learning method is
ultimately approximating \cite{vcl}. Mass survives only where \emph{every} task's posterior has
mass, so $q_t$ concentrates on the \textbf{intersection} of the solution
families. A network excellent at task 1 and hopeless at task 2 gets almost no
weight, however good task 1 alone would rate it.

\begin{remark}[We verified the intersection is non-empty]
Two networks trained from a \emph{shared} initialisation, one on task 1 and one
on task 2, each score $0.06$--$0.09$ on the other's task: they are specialists.
But the midpoint of the straight line between them scores $0.901$ and $0.889$.
The intersection is real and lies between the specialists.
\end{remark}

Splitting the last factor out of \eqref{eq:poe}, or applying Bayes directly,
gives the recursion:
\begin{equation}
q_{t+1}(\theta) \;\propto\; \underbrace{p(\Dt_{t+1} \mid \theta)}_{\Lh_{t+1}(\theta)} \; q_t(\theta).
\label{eq:seqbayes}
\end{equation}

\subsection{The real difficulty is representation}

Equation~\eqref{eq:seqbayes} says $q_t$ is a \emph{sufficient summary}:
everything the old data had to say about $\theta$ is in it, and the old datasets
are not needed. The difficulty is that we cannot store $q_t$.

It is a function on $\Rb^{89{,}610}$. Tabulating it is impossible. Writing it
symbolically as $\Lh_t \cdots \Lh_1 p(\theta)$ is exact but requires every
dataset to evaluate --- which is replay in different notation, and the data is
gone. That leaves lossy summaries, and the choice of summary \emph{is} the
method:

\begin{center}
\begin{tabular}{lll}
\toprule
method & representation of $q_t$ & what it discards \\
\midrule
EWC & point $\mu$ + diagonal $\Lambda$ & everything but a local quadratic \\
VCL & mean-field Gaussian & all correlation; re-projected each task \\
ours & $K$ particles + $(\mu, \Lambda)$ & whatever the $K$ do not cover \\
\bottomrule
\end{tabular}
\end{center}

\begin{remark}[Continual learning is recursive lossy compression]
Because \eqref{eq:seqbayes} feeds $q_t$ into the next step, each boundary's
compression error becomes the next boundary's starting error. Ten tasks is ten
rounds of re-encoding an already-degraded encoding. Forgetting, seen from the
inside, is accumulated compression error.
\end{remark}

\subsection{What we carry, and the asymmetry it creates}

We carry $K = 128$ particles $\Theta_t = \{\theta_t^{(k)}\}$ --- genuine
networks, trained then independently refined so they differ (spread $\approx
0.023$) --- together with a diagonal precision $\Lambda_t$.

When task $t+1$ arrives:
\begin{itemize}
\item Of $q_t$ we have \textbf{samples}: 128 arrays of 89,610 floats, each of
which can be loaded and run.
\item Of $q_{t+1}$ we have a \textbf{scoring procedure}: hand it any $\theta$
and \eqref{eq:seqbayes} returns a number proportional to $q_{t+1}(\theta)$.
\end{itemize}

Scoring is not generating. We can recognise a good network but cannot produce
one: $q_{t+1}$ concentrates in an unimaginably small region of
$\Rb^{89{,}610}$, and random weight vectors score at chance forever. Since
prediction requires actual networks to forward-pass, we must \emph{produce}
$K$ arrays distributed like $q_{t+1}$.

\begin{remark}[Why we cannot simply reweight]
Importance sampling is the textbook move: weight particle $k$ by
$\Lh_{t+1}(\theta^{(k)})$. It is exactly correct and it collapses. The
likelihood over $N$ examples is $\exp(-N\,\mathrm{CE})$, so with $N = 60{,}000$
a loss gap of $0.001$ nats between two particles produces a weight ratio of
$10^{26}$. All our particles are at chance on the new task and differ randomly,
so one takes all the mass. We measured the effective sample size falling to
$1.4$ out of $32$. Reweighting cannot move mass into a region the samples do not
occupy --- the particles must physically move.
\end{remark}

% ===========================================================================
\section{The moving target, and the chase}
\label{sec:chase}

This section describes what actually happens to the weights. Everything later
is machinery for doing this well.

\subsection{A sequence of targets}

We do not move from $q_t$ to $q_{t+1}$ in one jump. We define a sequence of
intermediate distributions and step through them. For a dial
$\beta \in [0,1]$,
\begin{equation}
\pi_\beta(\theta) \;\propto\; \Lh_{t+1}(\theta)^{\beta} \; q_t(\theta).
\label{eq:path}
\end{equation}

Read this in words. $\pi_\beta$ is \emph{``networks the old posterior likes,
reweighted by new-task performance, with the new task counted at strength
$\beta$.''}

\begin{itemize}
\item $\beta = 0$: the new task is ignored entirely. $\pi_0 = q_t$.
\item $\beta = 0.5$: a network is favoured if it is good on the old tasks and
somewhat good on the new one. Being excellent at the new task earns only half
the credit it eventually will.
\item $\beta = 1$: full Bayes. $\pi_1 = q_{t+1}$.
\end{itemize}

So \eqref{eq:path} is a smooth sequence of distributions beginning at what we
have and ending at what we want. Substituting \eqref{eq:seqbayes} shows it is
the geometric interpolation, $\pi_\beta \propto q_t^{1-\beta} q_{t+1}^{\beta}$ ---
the standard annealing path of sequential Monte Carlo and annealed flow
transport \cite{aft}.

\subsection{The picture: a target that moves, and particles that chase it}

Here is the mental model for the whole method.

\begin{center}
\emph{At every $\beta$ there is a target distribution $\pi_\beta$. Our 128
particles are trying to be samples of it. As $\beta$ advances the target moves,
and the particles chase.}
\end{center}

At $\beta = 0$ there is nothing to do: the particles already are samples of
$\pi_0 = q_t$. Advance $\beta$ slightly and the target shifts --- it now wants
networks tilted a little toward the new task. The particles are now slightly
wrong, and a short run of updates corrects them. Advance again, correct again.
After 25 such stages the target has become $q_{t+1}$, and if the particles kept
up, they are samples of it.

\subsection{What a single update does}

To chase $\pi_\beta$ we follow its gradient. Taking logs of \eqref{eq:path},
\begin{equation}
\log \pi_\beta(\theta) = \beta \log \Lh_{t+1}(\theta) + \log q_t(\theta) + \text{const},
\end{equation}
so
\begin{equation}
\nabla_\theta \log \pi_\beta(\theta)
=
\underbrace{-\,\beta\, \nabla_\theta L_{t+1}(\theta)}_{\text{new-task pull, scaled by }\beta}
\;\;
\underbrace{-\,\Lambda_t(\theta - \mu_t)}_{\text{pull back to the old solution}} .
\label{eq:twoforces}
\end{equation}

Every update applies exactly these two forces. Nothing else happens.

\begin{itemize}
\item The first is an ordinary gradient step on the new task, multiplied by
$\beta$. At stage 1 it is $4\%$ of a normal step; at stage 25 it is a full one.
\item The second is a spring pulling each weight back toward where it was, with
\textbf{a different stiffness for every weight}. Weight $j$ feels stiffness
$\Lambda_{t,j}$.
\end{itemize}

\begin{remark}[Where $\Lambda$ and $\mu$ come from]
The particles are the representation of $q_t$, but a finite set of points has no
gradient. To evaluate $\nabla \log q_t$ at an arbitrary $\theta$ we place a
Gaussian kernel around each particle,
$q_t(\theta) \approx \frac{1}{K}\sum_k \mathcal{N}(\theta; \theta^{(k)},
\Lambda^{-1})$, so $\mu$ is the particle's own previous position and $\Lambda$
is the width. Crucially $\Lambda$ is \emph{not} fitted to the particle spread
--- 128 points in $89{,}610$ dimensions span a degenerate subspace. It is the
accumulated Fisher, measured from data, which is the Laplace approximation's
curvature and says how sharply the log-posterior falls off in each direction.
\end{remark}

\subsection{Where the two forces balance}

The clearest way to see what the chase is doing is to ask where it would stop.
Setting the total force in \eqref{eq:twoforces} to zero:
\begin{equation}
\beta\,\nabla_\theta L_{t+1}(\theta) + \Lambda(\theta - \mu) = 0
\qquad\Longrightarrow\qquad
\boxed{\;\theta - \mu \;=\; -\,\beta\; \Lambda^{-1} \nabla_\theta L_{t+1}(\theta)\;}
\label{eq:equilibrium}
\end{equation}

Equation~\eqref{eq:equilibrium} is the mechanism of the method in one line. It
says how far from its old value each weight settles, and two features matter:

\paragraph{The displacement grows linearly in $\beta$.}
At $\beta$ small the particles sit close to where they started. As the dial
turns, the balance point slides steadily away from the anchor. \emph{The
equilibrium is the moving target, and it moves smoothly.}

\paragraph{Each coordinate is scaled by $1/\Lambda_j$.}
A weight with $\Lambda_j = 100$ barely moves even at $\beta = 1$: it was
important to the earlier tasks and the spring holds it. A weight with
$\Lambda_j \approx 0$ moves freely: nothing earlier depended on it, so the new
task may have it.

The spring term $-\Lambda(\theta-\mu)$ is the gradient of the EWC penalty
\cite{ewc}. The difference is that EWC \emph{postulates} that penalty, whereas
here it arrives as one of the two terms of $\nabla\log\pi_\beta$ --- accompanied
by the $\beta$ schedule, which EWC has no analogue of.

\textbf{This is how task 1 survives while task 2 is learned.} Not by moving a
little in every direction, but by moving a lot in the directions nobody needed
and hardly at all in the directions that mattered. A uniform spring could not do
this --- it would force one global compromise between plasticity and stability,
whereas \eqref{eq:equilibrium} resolves that trade-off separately for each of
the $89{,}610$ weights, using measured curvature.

\subsection{One stage, concretely}

With $K = 128$, $M = 25$ stages, $S$ updates per stage:

\begin{enumerate}
\item Set $\beta_m = m/25$. The equilibrium \eqref{eq:equilibrium} slides
slightly further from the anchor.
\item Record the two forces once, on a shared batch of $1024$ examples: the
new-task gradient $g_m$ and the spring force $a_m$. (These are what the learned
drift will be conditioned on; \S\ref{sec:implementation}.)
\item Run $S$ updates. Each draws 128 \emph{independent} minibatches of 64 ---
one per particle, so they stay distinct rather than becoming 128 copies of one
trajectory --- then applies the two forces of \eqref{eq:twoforces}.
\item Record the displacement $\Delta z_m$ that those $S$ updates produced.
\end{enumerate}

Measured magnitudes: the new-task step has RMS $8.9\times10^{-5}$, the spring
$1.7\times10^{-5}$. The spring is smaller on average but concentrated entirely
in the coordinates that matter, which is why it costs so little accuracy on the
new task while preserving the old.

% ===========================================================================
\section{The chase falls behind: why a bridge is needed}
\label{sec:lag}

\subsection{The ideal case}

Suppose the $S$ updates per stage were enough for the particles to fully settle
at $\pi_{\beta_m}$ before $\beta$ advanced. Then at every stage they would be
exact samples of the current target, and at $\beta = 1$ they would be exact
samples of $q_{t+1}$.

\begin{center}
\emph{If the dial turned infinitely slowly, we would be finished --- and there
would be no Schr\"odinger bridge in this paper.}
\end{center}

\subsection{What actually happens}

At finite speed the particles do not settle. The dial moves faster than they
relax, so at each stage they are still short of the current target, and they
enter the next stage already behind. The error accumulates. \textbf{At
$\beta = 1$ they land near $q_{t+1}$, not on it.}

This is not hypothetical. Increasing the compute per boundary makes the chase
tighter, and accuracy improves monotonically with no sign of saturation:

\begin{center}
\begin{tabular}{lcccc}
\toprule
updates per boundary & $2{,}500$ & $5{,}000$ & $10{,}000$ & $20{,}000$ \\
\midrule
accuracy & $0.9011$ & $0.9063$ & $0.9078$ & $0.9119$ \\
\bottomrule
\end{tabular}
\end{center}

Still climbing at $20{,}000$. \emph{That gap between where we are and the
infinitely-slow ideal is the lag.}

\subsection{And the schedule is what makes the chase work at all}

It is worth confirming that the tempering is doing something, rather than the
compute alone. Running the same two forces at full strength from the start ---
no dial, $\beta = 1$ throughout --- gives the opposite behaviour:

\begin{center}
\begin{tabular}{lcccc}
\toprule
updates per boundary & $2{,}500$ & $5{,}000$ & $10{,}000$ & $20{,}000$ \\
\midrule
with the dial & $0.9011$ & $0.9063$ & $0.9078$ & $\mathbf{0.9119}$ \\
without & $0.8890$ & $0.8834$ & $0.8783$ & $0.8789$ \\
\bottomrule
\end{tabular}
\end{center}

Without the dial, \textbf{more compute makes things worse} --- forgetting rises
from $0.073$ to $0.088$. The extra updates simply drive the particles further
into a solution for the new task alone. The two curves move in opposite
directions under identical budgets, so the gain is attributable to the path
rather than to the optimisation effort.

\subsection{So: what is the bridge for}

We want the accuracy of an infinitely slow chase at the cost of a fast one. The
Schr\"odinger bridge is precisely the object that provides it:

\begin{center}
\emph{the minimal correction to the chase that makes it land exactly on
$q_{t+1}$ despite running in finite time.}
\end{center}

Before we can say ``minimal'' we need to say minimal \emph{with respect to
what}, which is the next section.

% ===========================================================================
\section{Choosing the reference process}
\label{sec:reference}

A cost on transports is always measured relative to a \emph{reference process}:
a baseline notion of what motion is free. This choice is not a technicality. It
determines what the method considers expensive, which is to say what it
considers damage.

\subsection{Option 1: Brownian motion, and why it is wrong here}

The standard choice in the Schr\"odinger-bridge literature is
$\mathrm{d}\theta_s = \sigma\,\mathrm{d}W_s$ --- pure diffusion, no preferred
direction.

Under a Brownian reference, ``cheap'' means \emph{short in Euclidean distance},
and the reference's own bridges are Brownian bridges: straight lines plus noise.
The resulting transport would move each particle along the straight segment to
its destination, entirely blind to the loss landscape.

Two things go wrong, and both are fatal in weight space.

\paragraph{Euclidean distance does not measure damage.}
A large move in a direction no earlier task constrained is harmless. A tiny move
in a direction task 1 depended on is catastrophic. A Brownian reference prices
these identically, by length. It would happily destroy a critical weight so long
as it did so over a short distance.

\paragraph{Straight lines leave the manifold.}
The segment between two good networks passes through weight vectors that may
classify nothing. Nothing in a Brownian reference discourages this, because it
has no notion of which regions of weight space contain functioning networks.

\subsection{Option 2: annealed Langevin, which is the chase itself}

Instead we take as reference the process that is \emph{already doing what we
want} --- the chase of \S\ref{sec:chase}:
\begin{equation}
\mathrm{d}\theta_s \;=\; \nabla_\theta \log \pi_{\beta(s)}(\theta_s)\,\mathrm{d}s
\;+\; \sqrt{2}\,\mathrm{d}W_s,
\label{eq:reference}
\end{equation}
with the drift given by \eqref{eq:twoforces}: the two forces, nothing added.

At frozen $\beta$ this SDE has $\pi_\beta$ as its stationary distribution, so it
is exactly the dynamics whose equilibrium is \eqref{eq:equilibrium}. Advancing
$\beta$ makes it \emph{annealed} Langevin.

Three consequences, each of which fixes one of Brownian motion's problems.

\paragraph{``Free'' now means following the two forces.}
Motion the new-task gradient and the spring were already producing costs
nothing. Only \emph{additional} force is charged for. And additional force is by
definition movement the evidence does not justify --- which is the right thing
to penalise.

\paragraph{The metric becomes Fisher-weighted, automatically.}
Because the reference contains $-\Lambda(\theta-\mu)$, moving against that pull
requires extra force proportional to $\Lambda_j$. High-curvature directions are
expensive; flat ones are nearly free. \textbf{The reference encodes what counts
as damage, and this is where the per-weight geometry enters the cost.} We never
had to impose it separately.

\paragraph{The reference stays where networks work.}
Its trajectories are gradient descent trajectories on a tempered objective, so
they remain in regions of weight space where the loss is low --- not straight
lines through the void.

\subsection{And it very nearly solves the problem already}

The decisive property, restated from \S\ref{sec:lag}: run \eqref{eq:reference}
infinitely slowly and its marginal is $\pi_\beta$ at every $\beta$, so it lands
exactly on $q_{t+1}$. Run it at finite speed and it lags.

\begin{center}
\emph{$Q$ is a natural process that almost solves our problem, and fails only
because we run it in finite time.}
\end{center}

That is why it is the right reference. We are not importing a cost function by
analogy; we are asking how little we can perturb a process that was nearly
correct. And ``nearly correct'' means the correction is small, which
--- as \S\ref{sec:implementation} shows --- is what makes it learnable at all.

% ===========================================================================
\section{The Schr\"odinger bridge}
\label{sec:sb}

\subsection{Statement}

A \emph{path measure} $P$ is a distribution over whole trajectories
$(\theta_s)_{s\in[0,1]}$, not over points; $P_s$ is its marginal at time $s$ and
$P_{01}$ the joint law of its endpoints.

\begin{definition}[Schr\"odinger bridge \cite{leonard}]
\begin{equation}
P^\star = \operatorname*{arg\,min}_{P} \KL(P \| Q)
\qquad\text{subject to}\qquad P_0 = q_t, \quad P_1 = q_{t+1}.
\label{eq:sb}
\end{equation}
\end{definition}

\begin{remark}[The constraint is doing the work]
Without the constraint the answer is trivial: be $Q$, at zero cost. But $Q$ is
\emph{infeasible} --- it does not arrive. So $P^\star$ is the \emph{cheapest
feasible} process: the minimal correction to the chase that fixes where it
lands.
\end{remark}

For diffusions sharing a volatility, Girsanov makes the cost explicit:
\begin{equation}
\KL(P\|Q) \;=\; \tfrac{1}{2}\,\Eb_P \int_0^1 \big\lVert b^P_s(\theta_s) - b^Q_s(\theta_s) \big\rVert^2 \mathrm{d}s ,
\end{equation}
the integrated squared \emph{extra force} applied beyond the two natural ones.
Not distance travelled --- control effort.

\subsection{Why KL and not some other cost}

By Sanov's theorem, if a large population evolves under $Q$ and we
\emph{condition} on the atypical event that its terminal empirical distribution
is $q_{t+1}$, the conditional law converges to exactly the constrained
KL-minimiser. So \eqref{eq:sb} is not a modelling preference: it \emph{is} the
posterior over trajectories given the two observations. $Q$ is the prior over
motion, the marginals are the evidence.

\subsection{Two properties of the solution}

Neither is imposed; both follow, and together they make $P^\star$ computable.

\paragraph{It uses $Q$'s own routes (``reciprocal'').}
Split a path into endpoints and middle and apply the chain rule for KL:
\begin{equation}
\KL(P\|Q) = \KL\big(P(\text{ends})\|Q(\text{ends})\big)
+ \Eb\Big[\KL\big(P(\text{mid}\mid\text{ends}) \| Q(\text{mid}\mid\text{ends})\big)\Big].
\label{eq:decomp}
\end{equation}
Both terms are non-negative and the second involves only the middle, so it can
be set to \emph{exactly zero} by copying $Q$'s conditional middle --- which
costs nothing, since the constraints are on the ends. Once a particle's start
and destination are fixed, the cheapest route between them is however the
reference would have gone. \textbf{All the optimisation is therefore in the
coupling}: in deciding which particle becomes which solution.

\paragraph{It is Markov.}
If $Q$ is Markov --- and \eqref{eq:reference} is --- the constrained minimiser
is too. This is fortunate rather than incidental: a Markov process can be
simulated knowing only the current state, via a drift $v(\theta,s)$. Our entire
implementation, a network reading the current weights and the current forces,
is well posed only because of this.

\begin{theorem}[\cite{dsbm}]
\label{thm:char}
$P^\star$ is the \emph{unique} path measure that is both Markov and in the
reciprocal class $\Rcal(Q)$, with the prescribed marginals.
\end{theorem}

The two families genuinely differ. A member of $\Rcal(Q)$ picks an endpoint pair
first and then travels, so observing where a path has been reveals which pair it
is on and hence something about its future --- that is memory, and memory is not
Markov. The bridge is the one coupling for which the two coincide.

\subsection{The shape of the answer}

By the Doob $h$-transform,
\begin{equation}
v^\star(\theta,s) = \underbrace{\nabla_\theta \log \pi_{\beta(s)}(\theta)}_{\text{the chase}}
+ \underbrace{\nabla_\theta \log h(\theta,s)}_{\text{correction}},
\label{eq:htransform}
\end{equation}
where the correction says which way to lean \emph{now} so the arrival comes out
right. Computing $h$ exactly requires solving a PDE in $89{,}610$ dimensions, so
we learn the correction instead --- and \eqref{eq:htransform} dictates the
architecture: a correction on top of the reference, zero-initialised, so that an
untrained drift degrades to the plain chase rather than to noise.

% ===========================================================================
\section{Finding it: Iterative Markovian Fitting}
\label{sec:imf}

Theorem~\ref{thm:char} says $P^\star$ is the unique member of two families. We
cannot construct it, but we can enforce each family in turn.

\paragraph{Reciprocal projection --- keep the ends, restore $Q$'s middles.}
\begin{equation}
\proj_\Rcal(P) = P_{01} \otimes Q_{\mid 01}.
\end{equation}
That this is the KL projection is immediate from \eqref{eq:decomp}: for
$\Pi = \pi \otimes Q_{\mid01}$ the second term is independent of $\pi$, so the
minimiser is $\pi = P_{01}$. \emph{Result:} in $\Rcal(Q)$; \emph{but} the
particles carry destinations again, so Markovianity breaks.

\paragraph{Markovian projection --- erase the memory.}
\begin{equation}
v(\theta,s) = \Eb_P\!\left[\dot\theta_s \mid \theta_s = \theta\right],
\label{eq:markovproj}
\end{equation}
the average step taken from each state, rebuilt as a memoryless process. By
Girsanov the conditional expectation is the $L^2$-minimiser, so this is the KL
projection. \emph{Result:} Markov; \emph{but} the induced coupling is no longer
$Q$'s, so reciprocality breaks.

\paragraph{Alternate.}
Each projection repairs one property and damages the other, by less each round,
converging to the intersection. Both preserve the marginals --- $\proj_\Rcal$
leaves $P_{01}$ untouched and $\proj_\Mcal$ reproduces every time marginal ---
so the constraints hold at every iterate.

\begin{remark}[One pass is flow matching; iterating is the bridge]
\eqref{eq:markovproj} is a regression. Applying it once, to whatever coupling
you began with, gives flow matching \cite{flowmatch} --- and that is where most
implementations stop. But that process's coupling is whatever your initial guess
produced, with no reason to be the reciprocal one. The re-coupling step is the
difference between a curve fit and a bridge.
\end{remark}

% ===========================================================================
\section{Implementation}
\label{sec:implementation}

\subsection{The prior operator, solved exactly}

Rather than a gradient step on $\log q_t$ we apply its exact proximal operator:
\begin{align}
\proj(\theta) &= \operatorname*{arg\,min}_x \frac{\lVert x-\theta\rVert^2}{2\alpha\eta}
+ \tfrac12 (x-\mu)^\top \Lambda (x-\mu)
\\
\Longrightarrow\quad x_j &= \frac{\theta_j + \alpha\eta\Lambda_j \mu_j}{1+\alpha\eta\Lambda_j} .
\label{eq:prox}
\end{align}
A weighted average between where the weight is and where it was.

\begin{remark}[Operator, not penalty]
Adding $\tfrac{\lambda}{2}\lVert\theta-\mu\rVert^2_\Lambda$ to the loss routes
the prior gradient through Adam, which divides each coordinate by its own
gradient RMS --- cancelling exactly the curvature weighting the penalty exists
to impose. Applied outside the preconditioner, the weighting survives.
\end{remark}

\subsection{Tokenisation}

A whole network is one sample in $89{,}610$ dimensions and we hold $128$ of
them: $0.0014$ samples per dimension, hopeless for any regression. Viewing a
network instead as $210$ per-neuron weight vectors (layer $\ell$ contributes
$n_\ell$ tokens of dimension $d_\ell+1$) gives $26{,}880$ tokens in total, but
tokens from different layers live in different spaces and cannot be pooled, so
the ratio must be read per layer: $100\times128 = 12{,}800$ samples in $785$
dimensions ($16.3$ per dimension) for the first layer, $12{,}800$ in $101$
($126.7$) for the second, and $10\times128 = 1{,}280$ in $101$ ($12.7$) for the
head. The worst layer is still four orders of magnitude better off than the
whole-network view, and the map is affine and invertible.

\subsection{The learned drift}

The Markovian projection \eqref{eq:markovproj} becomes
\begin{equation}
\min_\phi \Eb \left\lVert d_\phi(z_m, g_m, a_m, \beta_m) - \Delta z_m \right\rVert^2 ,
\end{equation}
with $\Delta z_m$ the displacement the reference updates produced over stage
$m$. Conditioning on $g_m$ and $a_m$ --- the two terms of \eqref{eq:twoforces}
--- is what makes the fitted drift a \emph{rule} (how far to move given this
much new-task pressure against this much spring resistance at this $\beta$)
rather than a route memorised for one task, and hence transferable to later
boundaries.

\subsection{IMF across the stream}

Each boundary: traverse and record; fit $d_\phi$ forward and $d_\psi$ backward
$(\proj_\Mcal)$; discard the recording so the next boundary regenerates the
coupling under the drift just fitted $(\proj_\Rcal)$. The task stream supplies
the iterations. Convergence is monitored by the cycle residual
\begin{equation}
\mathcal{C} = \frac{\Eb\lVert d_\phi(z) + d_\psi(z + d_\phi(z))\rVert^2}{\Eb\lVert d_\phi(z)\rVert^2},
\end{equation}
zero exactly on $\Mcal \cap \Rcal(Q)$.

\subsection{One approximation, stated plainly}

We run the deterministic flow: the $\sqrt2\,\mathrm{d}W_s$ of
\eqref{eq:reference} is off. The path $\pi_\beta$ is unchanged but particles are
transported rather than sampled, so diversity comes from initialisation and
per-particle minibatches. This is why anchoring each particle to its own
previous position rather than to the population mean mattered empirically
(spread $0.029$ versus $0.001$).

% ===========================================================================
\section{What we measured}
\label{sec:numbers}

\begin{center}
\begin{tabular}{lccc}
\toprule
& forward expl.\ var. & backward expl.\ var. & cycle $\mathcal{C}$ \\
\midrule
after 1 round (bridge matching) & $0.286$ & $0.168$ & $0.908$ \\
after 2 rounds (IMF) & $0.624$ & $0.666$ & $\mathbf{0.114}$ \\
\bottomrule
\end{tabular}
\end{center}

An eightfold reduction in the cycle residual: the iterations do move the fitted
pair toward $\Mcal \cap \Rcal(Q)$.

\begin{center}
\begin{tabular}{lcc}
\toprule
traversal & updates per boundary & average accuracy \\
\midrule
reference chase & $2500$ & $0.9099$ \\
learned drift, IMF $=2$ & $1250$ & $0.9086$ \\
learned drift, IMF $=1$ & $625$ & $0.9066$ \\
learned drift, IMF $=2$ & $625$ & $0.9039$ \\
\bottomrule
\end{tabular}
\end{center}

\begin{remark}[The dissociation]
Solving the bridge markedly better ($\mathcal{C}: 0.908 \to 0.114$) does not
improve accuracy. Transport quality was not the binding constraint: accuracy
comes from the path together with the Fisher geometry and the ensemble, and once
on the path, traversing it more exactly buys \emph{compute} rather than
retention. This is consistent with Kessler et al.\ \cite{kessler}, who find that
even exact sequential Bayesian inference fails to prevent forgetting, attributing
it to model misspecification rather than inference quality.
\end{remark}

\begin{thebibliography}{9}

\bibitem{leonard}
C.~L\'eonard, ``A survey of the Schr\"odinger problem and some of its connections
with optimal transport,'' \emph{Discrete Contin.\ Dyn.\ Syst.}, vol.~34, no.~4,
pp.~1533--1574, 2014.

\bibitem{dsbm}
Y.~Shi, V.~De~Bortoli, A.~Campbell, and A.~Doucet, ``Diffusion Schr\"odinger
bridge matching,'' in \emph{Proc.\ NeurIPS}, 2023.

\bibitem{flowmatch}
Y.~Lipman, R.~T.~Q. Chen, H.~Ben-Hamu, M.~Nickel, and M.~Le, ``Flow matching for
generative modeling,'' in \emph{Proc.\ ICLR}, 2023.

\bibitem{aft}
M.~Arbel, A.~G.~D.~G. Matthews, and A.~Doucet, ``Annealed flow transport Monte
Carlo,'' in \emph{Proc.\ ICML}, 2021.

\bibitem{kessler}
S.~Kessler, A.~Cobb, T.~G.~J. Rudner, S.~Zohren, and S.~J. Roberts, ``On
sequential Bayesian inference for continual learning,'' \emph{Entropy}, vol.~25,
no.~6, 2023.

\bibitem{vcl}
C.~V. Nguyen, Y.~Li, T.~D. Bui, and R.~E. Turner, ``Variational continual
learning,'' in \emph{Proc.\ ICLR}, 2018.

\bibitem{ewc}
J.~Kirkpatrick et al., ``Overcoming catastrophic forgetting in neural networks,''
\emph{PNAS}, vol.~114, no.~13, pp.~3521--3526, 2017.

\end{thebibliography}

\end{document}
